% Stand 27. Juli 2026
%\section*{Zusammenfassung}
\label{KapitelZusammenfassung}
Gegenstand dieser Arbeit ist die räumliche Analyse von
Feinstaubmessungen im Großraum München unter Verwendung
geostatistischer Methoden. \\


Im theoretischen Kapitel 2 werden zunächst die Grundlagen der Geostatistik
bereitgestellt. Semivarianz und Variogramm werden als zentrale Werkzeuge 
der räumlichen Modellierung eingeführt, gefolgt von einer Übersicht 
gängiger Kovarianz-Modelle sowie der Fachbegriffe Nugget, Sill und Range.
Anschließend werden die Kriging-Gleichungen hergeleitet, die sich als 
Hierarchie zunehmend allgemeinerer Annahmen an den Erwartungswert 
verstehen lassen. 
Das Kapitel schließt mit einer Darstellung der geometrischen und 
zonalen Anisotropie.\\

Das darauf folgende Kapitel 3 beschreibt Herkunft und Struktur der verwendeten
Daten, die mittels kostengünstiger Sensoren gemessen und über das 
Bürger-Netzwerk \textsc{SensorCommunity} zum Download bereitgestellt werden. 
Nach einer mehrstufigen Datenbereinigung stehen für drei Monate bereinigte,
monatlich gemittelte Messreihen von rund 80 bis 95 Sensoren zur Verfügung. \\

Kapitel~4 widmet sich der zentralen Fragestellung, ob die Daten Hinweise auf
Anisotropie liefern. Die Schätzung der direktionalen Variogramme zeigt, dass die maximale 
paarweise Distanz der Sensoren sehr wahrscheinlich  
unterhalb der tat\-säch\-li\-chen Reichweite der räumlichen Korrelation liegt. Als Folge lässt sich der 
Range-Parameter nicht zuverlässig schätzen, so dass
die Standardmethode der Variogramm-Karten keine 
schlüssigen Ergebnisse liefert. Allerdings deuten 
die Richtungs\-ab\-hängig\-keit des Sill sowie eine 
explorative Methode mit konstruierten 
anisotropen Kovarianz-Modellen konsistent auf eine bevorzugte Richtung hin. 
Diese Befunde sind als Indizien für eine mögliche
geometrische Anisotropie zu werten. Ein statistisch
belastbarer Beweis gelingt aufgrund der begrenzten räumlichen 
Abdeckung und der nicht gegebenen stochastischen Unabhängigkeit 
jedoch nicht. \\

Kapitel~5 fasst zusammen, dass die Datenlage keine wissenschaftlich haltbare
Entscheidung zwischen geometrischer und zonaler Anisotropie erlaubt, anisotrope Modelle
aber eine sinnvolle Erweiterung isotroper Ansätze darstellen können. Für eine 
verlässliche Bestätigung sind jedoch umfangreichere Datensätze 
über ein größeres Untersuchungsgebiet bei vergleichbarer 
Sensordichte erforderlich. 

